% !TeX root = ../papers/nc-p3.tex
\section{一般方向にも残る交換子復元の階数降下}
\label{sec:marked-persistent-bracket-defect}
\begin{lemma}\label{lem:nine-symmetric-words}
跡零三行列 $X,Y$ で $\det[X,Y]\ne0$ とする。$S=XY+YX$ と置くと
\[
                I,X,Y,X^2,Y^2,S,XYX,YXY,S^2
\]
は全三行列の基底である。
\end{lemma}
\begin{proof}
行優先で九列を並べた行列式は $-\det[X,Y]^3$ であることを示す。
まず $X$ を相異なる三固有値 $x_i$ の対角行列とする。$\det[X,Y]$ はVandermondeの
非零倍と $f=y_{12}y_{23}y_{31}-y_{13}y_{32}y_{21}$ の積である。
$f$ は既約で、$f=0$ の一般点では全非対角成分が非零となり、
$x_i$ の基底で対角な可逆計量に関して $Y$ も自己随伴になる。
九語は全て反転で不変なので計量に関して自己随伴となり、列階数は6以下である。
従って高さ1の $f$ で局所化して行列を掃き出すと行列式は $f^3$ で割れる。
$Y$ に関する全次数は9なので、商は $Y$ に依存しない。

$Y=E_{12}+E_{23}+E_{31}$ で商を計算する。対角、巡回一段、巡回二段の各三列の係数は、
巡回添字を使い、順に $(1,x_i,x_i^2)$, $(1,x_i+x_{i+1},x_ix_{i+1})$,
$(1,x_{i+1},(x_i+x_{i+1})(x_{i+1}+x_{i+2}))$ である。
跡零性から最後の二群の二次係数はそれぞれ $x_{i+2}^2-\operatorname{tr}(X^2)/2$ と
$x_{i+1}^2-\operatorname{tr}(X^2)/2$ になる。従って三つとも同じVandermondeの符号倍で、
$(x_1,x_2,x_3)=(1,2,-3)$ で順序の符号を決めると表示の恒等式を得る。
共役の全行列空間上の行列式は1であり、対角化可能な稠密開集合から全行列に延びる。
\end{proof}

\begin{theorem}\label{thm:persistent-three-bracket-defect}
既約な跡零三平面 $\mathcal Q$ の交換子写像 $B$ が三独立で、像が可逆元を含むとする。
一般の $X\in\mathcal Q$ でも前節の復元の階数が3に留まる必要十分条件は、
$U=k^2\oplus k$ の適当な基底で
\[
 \mathcal Q=\left\{\begin{pmatrix}aI_2&w\\ w^t S&-2a\end{pmatrix}
               :a\in k,\ w\in k^2\right\},\qquad S\in\mathrm{GL}_2,\quad S\ne S^t
\]
となることである。他の全ての三平面では一般方向の線型復元は階数4である。
表示の場合、判別三次式は非零定数を除いて一つの一次式の三乗であり、
$\mathbf P(kI+\mathcal Q)$ の階数一切断は平方零の平面三重点である。
\end{theorem}
\begin{proof}[二次積と特異方向]
$X$ の鉛筆 $N=[X,\mathcal Q]$ が可逆元を含む開集合で考える。
例外方向 $K=[X,Y]$ と補基底 $Z$ に対し $\{X,K\}=\alpha K+\beta [X,Z]$ とする。
$X$ は巡回であり、$[X,\{X,Y\}]=\{X,K\}$ だから
$\{X,Y\}-\alpha Y-\beta Z$ は $X$ の二次以下の多項式である。
$Y$ に $X$ の適当な倍を足した $Y'$ を使えば $\{X,Y'\}\in kI+\mathcal Q$。
従って $XY'$ は積を対称化する写像
$\mathrm{Sym}^2\mathcal Q\to\operatorname{End}(U)/(kI+\mathcal Q)$ の核に入る可約二次式である。

可逆な $K$ を持つ例外方向が一つでもあれば $\beta\ne0$ であり、
$\mathcal Q=\langle X,Y',\{X,Y'\}_0\rangle$ である。
このとき先の対称積の核はちょうど $k(XY')$ である。
実際補題の九語で $I,X,Y',\{X,Y'\}$ に
$X^2,Y'^2,\{X,\{X,Y'\}\},\{Y',\{X,Y'\}\},\{X,Y'\}^2$ を加えるとやはり基底となる。
$\{X,\{X,Y'\}\}$ は $XY'X$ に低次の $I,X,Y'$ の線型結合を加えたもので、
$X$ と $Y'$ を交換しても同じである（跡零三行列のCayley--Hamiltonの偏極）。
よって例外になり得る方向は二次式の二因子 $[X],[Y']$ だけで、一般方向は階数4である。

一般に階数3が残るとすれば、一般の鉛筆の例外は特異な $K$ である。
前節よりこれは判別三次曲線との重交点である。一般の直線は被約台の滑らかな点で横断するので、
三次式には唯一の重複成分である直線がある。この直線を $\mathbf P([T,\mathcal Q])$ とする
固定方向 $[T]\in\mathbf P(\mathcal Q)$ があり、一般の $X$ で $K=[X,T]$ としてよい。
ここで $\beta\ne0$ なら $\mathcal Q$ は $X,T$ およびその対称積と $X^2$ から生成される。
$\det[X,T]=0$ のため二行列には非零の共通対称形式がある
（定理\ref{thm:rank-eight-finite-symmetry}の六式の行列式恒等式）。
その核は両行列に不変で従って $\mathcal Q$ に不変、よって形式は可逆である。
全 $\mathcal Q$ がその形式に関して自己随伴になり、全交換子が特異になって矛盾する。
ゆえに全 $X$ で多項式恒等式
\[
                    [X^2,T]=\lambda(X)[X,T]
\]
が成立する。倍率が一次式であるのは $[X,T]$ の成分が二つ以上の独立な一次式を張り、
共通因子を持たないからである。$X=T+tY$ として一次項を比べれば
$[Y,T^2-\lambda(T)T]=0$。既約性により $T^2-\lambda(T)T$ はスカラーであり、
$T$ は非スカラー非巡回なので $P_0$ の非零倍、$P$ は階数1である。
\end{proof}

\begin{proof}[階数一方向に沿うブロック]
スカラーを引いても不変部分空間は変わらない。$P=E_{13}$ なら
$X-\lambda(X)I/2$ の平方が $P$ と交換する。
その九成分を前節までと同じ順に $a,b,c,d,e,f,g,h,i$ と書くと
$g(a+i)+hd=0$, $d(a+e)+fg=0$, $h(e+i)+gb=0$ である。
$g=0$ なら $d,h$ の一方が恒等零で、直線または平面が不変となる。
$g\ne0$ なら $g$ が $hd$ を割るので、必要なら反対角転置で双対に移って $d$ が $g$ の倍とできる。
$P$ を保つ旗の基底変更で $d=0$ とすると $f=0$ となり、$\langle e_1,e_3\rangle$ が不変である。

半単純 $P$ では $X=\left(\begin{smallmatrix}A&p\\q^t&b\end{smallmatrix}\right)$ とし、
$l=\lambda(X)-b$ とすると $A p=l p$, $q^t A=l q^t$。
$p,q$ の一方が恒等零なら可約なので、$A-lI$ の全元の階数は1以下である。
この空間が非零で共通像直線を持てば、$q$ は常にこの像を消し、$A$ は像直線を保つため
共通不変直線がある。共通核直線なら $p$ がそこに入り、最後の軸と核直線の張る平面が共通不変である。
よって $A=lI$ が恒等的に成立する。$P_0$ を使って他の二基底は両対角が零にできる。

この二元の交換子 $L$ はブロック対角であり、他の二交換子はブロック非対角である。
$zL+C$ の行列式は非対角部分 $C$ の符号に不変で、$z^3$ と $z$ の項だけを持つ。
$z=0$ は先に得た重複直線なので $\det L\ne0$ である。
従って左上ブロック $p_1 q_2^t-p_2 q_1^t$ も可逆で、$p_1,p_2$ と $q_1,q_2$ はそれぞれ独立。
これで表示の $w^t S$ が得られ、$L$ の最後の成分が非零のため $S$ は非対称である。
逆に表示では $P=E_{33}$ とすると $X^2-(a-2a)X$ はブロック対角で、
$\{X,[X,P]\}$ は $[X,P]$ の倍である。二つの座標方向の交換子は可逆で
（左上は可逆交代行列と $S$ の積、右下は $S_{12}-S_{21}$ の符号倍）、
$P_0$ との二交換子も独立である。表示の三平面は既約である。
実際 $P$ により不変空間は二ブロックの部分空間の和で、二組の非対角成分はいずれも二次元を張る。
よって一般方向に階数3が生じる。上の行列式の議論から三次式は三重直線となる。

最後に $kI+\mathcal Q$ で階数1なら幾何点は $P$ のみである。
右下を1とする近傍では式は $cI_2=w w^t S$ で、対称二行列から
$\mathrm{Mat}_2/kI$ への $M\mapsto MS$ は $S^{-1}$ が非対称のため同型である。
従って $w$ の三つの二次単項式と $c$ が零になり、主張の長さ3局所環を得る。
\end{proof}

\begin{corollary}\label{cor:no-persistent-defect-finite}
単純印を持つ片側対称有限窓で $q=3$, $g>3$、三交換子が独立なら、
復元の線型系は一般方向で必ず階数4である。対称三平面の和の次元4,5の両方を含む。
\end{corollary}
\begin{proof}
前定理の模型とし、$Q_s$ が $P_0$ の倍となる方向を $s$ とする。
$W=k^2$ を上のブロック、$A$ を最後の軸を消す二次元の行空間とする。
$G'=G_A$, $F'=F_W$ はそれぞれ二平面である。
$\det B(v)$ の三重直線を $\ell=0$ とすれば、その上の全行列は非対角ブロックだけを持つ。
ここで $B(v)$ は $v=x\times y$ に対応する指定値 $[Q_x,Q_y]$。
この直線上では列・行の各ブロックが二変数の同型となり、全点で階数2、左核は $A$ 内にある。
ブロックの一成分が単元の近傍で掃き出せばこの核の包含はスキーム上でも成立する。

$Z_G=\{[v]: vv^t\in kI+G\}\subset\mathbf P^2$ とする。
単独の四空間の正次元階数一切断は母線だけで、対称行列の二次Veroneseには母線がないので
$Z_G$ は有限である。二つの二次式の完全交叉ゆえ長さ4である。
$Z_G$ 上で $T_v=vv^t-(v^tv)I/3$ は $G$ の非零元であり、外積行列 $[v]_\times$ と交換する。
従って全 $F_c$ に対し $\operatorname{tr}([v]_\times[T_v,F_c])=0$。
$G$ 基底での $T_v$ の係数が局所的に単元成分を持つ左核を与えるので、
$\det B(v)$、従って $\ell^3$ が $Z_G$ 上で恒等零になる。
直線でさらに切ると左核は $A$ にあるため、$Z_G\cap V(\ell)$ 全体で
$vv^t\in kI+G'$ となる。$\ell$ の冪での三段のフィルトレーションから
$4\le3\operatorname{length}(Z_G\cap V(\ell))$ なので、後の切断は長さ2以上である。

これで $G'$ が可約と分かる。実際その平面 $\mathbf P(kI+G')$ には
Veronese上の長さ2部分の包が入る。相異なる二点なら対応する行列の像は二列の張る二平面にあり、
接線なら $uu^t,u v^t+v u^t$ の像はやはり二平面にある。
どちらもその包には単位行列がないから、これに $I$ を加えると $kI+G'$ 全体となり、
像の二平面を保つ。対称性よりその直交直線 $kn$ も保たれる。

$A,W$ の変数の交叉交換子は模型の左上ブロックから
$[G_a,F_c]=m(a,c)[s]_\times$ と書け、二次元の対合 $m$ は完全である。
実際 $[Q_s,Q_t]$ の左上は全て零、残る二つの非対角生成元の交換子の左上は可逆である。
従って $a$ に $m(a,c)=0$ の方向を対応させる二空間の線型同型 $a\mapsto c(a)$ がある。
一般の $G_a$ は巡回である。全二平面が非巡回なら、
命題\ref{prop:marked-three-commutators-reconstruction}の判別式の議論により
一次スカラーを引いた対称行列が全て階数1以下となる。しかしそのような対称線型空間は一次元以下である。
$F_{c(a)}$ は $G_a$ と交換するためその多項式で、$kn$ を保つ。
稠密性で全 $F'$ が $kn$ を保ち、従って全ての交叉交換子は $n$ を消す。
非零の $[s]_\times$ の核は $ks$ なので $n$ は $s$ に比例する。
$Q_s$ は $W$ 上でスカラーであり、全 $c\in W$ で $F_c s$ は $s$ の倍である。
定理\ref{thm:marked-commuting-intersection}の最後の辺 $12$ の六式を使うと、
$\mathbf P(W)$ の全直線が同じ隣接点の上に入り、有限性に反する。
\end{proof}
